% TikZ 图：匀变速直线运动的 v-t 图像
% Day 03：加速度
% 使用：% TikZ 图：匀变速直线运动的 v-t 图像
% Day 03：加速度
% 使用：% TikZ 图：匀变速直线运动的 v-t 图像
% Day 03：加速度
% 使用：% TikZ 图：匀变速直线运动的 v-t 图像
% Day 03：加速度
% 使用：\input{figures/tikz/vt-graph.tex}
\begin{tikzpicture}[scale=1,>=Latex]
  % 坐标轴
  \draw[->] (-0.3,0) -- (5.5,0) node[right] {$t$};
  \draw[->] (0,-0.3) -- (0,4.8) node[above] {$v$};
  % 标注原点
  \node[below left] at (0,0) {$O$};
  % 画 v-t 线（黑色粗实线）
  \draw[very thick] (0,1) -- (4,4);
  % v0 箭头
  \draw[->,thick] (0,0) -- (0,1);
  \node[left] at (0,0.5) {$v_0$};
  % t 时刻虚线到顶点，标 v
  \draw[dashed] (4,0) -- (4,4);
  \draw[->,thick] (4,0) -- (4,4);
  \node[right] at (4,2) {$v$};
  \node[below] at (4,0) {$t$};
  % 梯形面积标注
  \fill[black!15] (0,0) -- (0,1) -- (4,4) -- (4,0) -- cycle;
  \node at (2.5,0.8) {$x = \frac{v_0+v}{2}t$};
  % Δv / Δt 直角三角形：三边用虚线，顶点严格在直线上
  \draw[dashed,thick] (1,1.75) -- (3.5,1.75);  % Δt 水平边
  \draw[dashed,thick] (3.5,1.75) -- (3.5,3.625);  % Δv 竖直边
  \draw[dashed,thick] (1,1.75) -- (1.5,2.125);  % 斜边一段对齐直线
  % 标注 Δt 和 Δv
  \node[below] at (2.25,1.75) {$\Delta t$};
  \node[right] at (3.5,2.7) {$\Delta v$};
  % 斜率公式标注
  \node[above right] at (1,1.75) {$a = \frac{\Delta v}{\Delta t}$};
\end{tikzpicture}

\begin{tikzpicture}[scale=1,>=Latex]
  % 坐标轴
  \draw[->] (-0.3,0) -- (5.5,0) node[right] {$t$};
  \draw[->] (0,-0.3) -- (0,4.8) node[above] {$v$};
  % 标注原点
  \node[below left] at (0,0) {$O$};
  % 画 v-t 线（黑色粗实线）
  \draw[very thick] (0,1) -- (4,4);
  % v0 箭头
  \draw[->,thick] (0,0) -- (0,1);
  \node[left] at (0,0.5) {$v_0$};
  % t 时刻虚线到顶点，标 v
  \draw[dashed] (4,0) -- (4,4);
  \draw[->,thick] (4,0) -- (4,4);
  \node[right] at (4,2) {$v$};
  \node[below] at (4,0) {$t$};
  % 梯形面积标注
  \fill[black!15] (0,0) -- (0,1) -- (4,4) -- (4,0) -- cycle;
  \node at (2.5,0.8) {$x = \frac{v_0+v}{2}t$};
  % Δv / Δt 直角三角形：三边用虚线，顶点严格在直线上
  \draw[dashed,thick] (1,1.75) -- (3.5,1.75);  % Δt 水平边
  \draw[dashed,thick] (3.5,1.75) -- (3.5,3.625);  % Δv 竖直边
  \draw[dashed,thick] (1,1.75) -- (1.5,2.125);  % 斜边一段对齐直线
  % 标注 Δt 和 Δv
  \node[below] at (2.25,1.75) {$\Delta t$};
  \node[right] at (3.5,2.7) {$\Delta v$};
  % 斜率公式标注
  \node[above right] at (1,1.75) {$a = \frac{\Delta v}{\Delta t}$};
\end{tikzpicture}

\begin{tikzpicture}[scale=1,>=Latex]
  % 坐标轴
  \draw[->] (-0.3,0) -- (5.5,0) node[right] {$t$};
  \draw[->] (0,-0.3) -- (0,4.8) node[above] {$v$};
  % 标注原点
  \node[below left] at (0,0) {$O$};
  % 画 v-t 线（黑色粗实线）
  \draw[very thick] (0,1) -- (4,4);
  % v0 箭头
  \draw[->,thick] (0,0) -- (0,1);
  \node[left] at (0,0.5) {$v_0$};
  % t 时刻虚线到顶点，标 v
  \draw[dashed] (4,0) -- (4,4);
  \draw[->,thick] (4,0) -- (4,4);
  \node[right] at (4,2) {$v$};
  \node[below] at (4,0) {$t$};
  % 梯形面积标注
  \fill[black!15] (0,0) -- (0,1) -- (4,4) -- (4,0) -- cycle;
  \node at (2.5,0.8) {$x = \frac{v_0+v}{2}t$};
  % Δv / Δt 直角三角形：三边用虚线，顶点严格在直线上
  \draw[dashed,thick] (1,1.75) -- (3.5,1.75);  % Δt 水平边
  \draw[dashed,thick] (3.5,1.75) -- (3.5,3.625);  % Δv 竖直边
  \draw[dashed,thick] (1,1.75) -- (1.5,2.125);  % 斜边一段对齐直线
  % 标注 Δt 和 Δv
  \node[below] at (2.25,1.75) {$\Delta t$};
  \node[right] at (3.5,2.7) {$\Delta v$};
  % 斜率公式标注
  \node[above right] at (1,1.75) {$a = \frac{\Delta v}{\Delta t}$};
\end{tikzpicture}

\begin{tikzpicture}[scale=1,>=Latex]
  % 坐标轴
  \draw[->] (-0.3,0) -- (5.5,0) node[right] {$t$};
  \draw[->] (0,-0.3) -- (0,4.8) node[above] {$v$};
  % 标注原点
  \node[below left] at (0,0) {$O$};
  % 画 v-t 线（黑色粗实线）
  \draw[very thick] (0,1) -- (4,4);
  % v0 箭头
  \draw[->,thick] (0,0) -- (0,1);
  \node[left] at (0,0.5) {$v_0$};
  % t 时刻虚线到顶点，标 v
  \draw[dashed] (4,0) -- (4,4);
  \draw[->,thick] (4,0) -- (4,4);
  \node[right] at (4,2) {$v$};
  \node[below] at (4,0) {$t$};
  % 梯形面积标注
  \fill[black!15] (0,0) -- (0,1) -- (4,4) -- (4,0) -- cycle;
  \node at (2.5,0.8) {$x = \frac{v_0+v}{2}t$};
  % Δv / Δt 直角三角形：三边用虚线，顶点严格在直线上
  \draw[dashed,thick] (1,1.75) -- (3.5,1.75);  % Δt 水平边
  \draw[dashed,thick] (3.5,1.75) -- (3.5,3.625);  % Δv 竖直边
  \draw[dashed,thick] (1,1.75) -- (1.5,2.125);  % 斜边一段对齐直线
  % 标注 Δt 和 Δv
  \node[below] at (2.25,1.75) {$\Delta t$};
  \node[right] at (3.5,2.7) {$\Delta v$};
  % 斜率公式标注
  \node[above right] at (1,1.75) {$a = \frac{\Delta v}{\Delta t}$};
\end{tikzpicture}
